% =====================================================================
\chapter{Что такое EdgeTX}
% =====================================================================
\chapterintro
  {откуда взялась EdgeTX, как устроен путь сигнала внутри пульта, как устроены меню и как
   в них перемещаться кнопками TX12 и Boxer}
  {включённый пульт}

\section{Немного истории}

В начале 2010-х энтузиасты написали открытую прошивку \textbf{OpenTX} для пультов FrSky Taranis.
Её гибкость сделала открытые пульты стандартом в FPV и авиамоделизме. В 2021 году часть
разработчиков отделилась и создала форк \textbf{EdgeTX}, чтобы быстрее добавлять поддержку
новых пультов, цветные экраны с виджетами, сенсорный ввод и ExpressLRS. Развитие OpenTX
практически остановилось, и сегодня почти все пульты RadioMaster, Jumper, Flysky (серия PL/EL)
и других производителей выпускаются сразу с EdgeTX.

\begin{itemize}
  \item EdgeTX --- \textbf{бесплатная} и \textbf{открытая} (GPLv2) прошивка;
  \item исходный код и выпуски: \code{github.com/EdgeTX/edgetx};
  \item документация: \code{edgetx.org} и \code{manual.edgetx.org};
  \item инструменты: \textbf{EdgeTX Buddy} (браузерная утилита для прошивки и подготовки SD-карты)
        и \textbf{EdgeTX Companion} (редактор моделей и симулятор на компьютере).
\end{itemize}

Нумерация версий --- \code{2.X.Y}: \code{X} меняется при появлении новых возможностей,
\code{Y} --- при исправлениях ошибок. Модели, сохранённые в старой версии, при обновлении
конвертируются автоматически; обратная совместимость (откат на старую версию с новыми
моделями) \emph{не гарантируется} --- поэтому перед обновлением всегда делается резервная копия.

\section{Путь сигнала: главная идея EdgeTX}

Чтобы уверенно настраивать EdgeTX, нужно держать в голове одну схему. Всё, что делает пульт,
--- это ежесекундно сотни раз прогоняет положения органов управления через конвейер:

\begin{figure}[H]
\centering
\begin{tikzpicture}[node distance=5mm and 7mm,every node/.style={font=\sffamily\small}]
  \node[blk,fill=blksw,minimum width=22mm] (src) {Источники\\\scriptsize стики, крутилки,\\\scriptsize переключатели};
  \node[blk,fill=blkin,right=of src,minimum width=20mm] (inp) {Inputs\\\scriptsize (входы)};
  \node[blk,fill=blkmix,right=of inp,minimum width=20mm] (mix) {Mixes\\\scriptsize (миксеры)};
  \node[blk,fill=blkout,right=of mix,minimum width=20mm] (out) {Outputs\\\scriptsize (выходы)};
  \node[blk,fill=blkrf,right=of out,minimum width=18mm] (rf) {Радио-\\модуль};
  \draw[arr] (src) -- (inp);
  \draw[arr] (inp) -- node[above,lbl]{I1…I32} (mix);
  \draw[arr] (mix) -- node[above,lbl]{CH1…CH32} (out);
  \draw[arr] (out) -- (rf);
  \node[below=4mm of inp,lbl,align=center,text width=27mm] {расходы, экспо,\\трим, пилотский\\«вкус»};
  \node[below=4mm of mix,lbl,align=center,text width=27mm] {какой вход идёт\\в какой канал,\\смешивание};
  \node[below=4mm of out,lbl,align=center,text width=27mm] {механика:\\реверс, пределы,\\сабтрим};
  \node[blk,fill=blksw,above=8mm of mix,minimum width=78mm] (ls) {Flight modes, Logical switches, Special functions, GVars, Curves};
  \draw[arr,dashed] (ls) -- (inp);
  \draw[arr,dashed] (ls) -- (mix);
  \draw[arr,dashed] (ls) -- (out);
\end{tikzpicture}
\caption{Конвейер обработки сигнала в EdgeTX}
\label{fig:pipeline}
\end{figure}

\begin{description}[font=\sffamily\bfseries\color{etx}]
  \item[Источники (sources)] всё, у чего есть значение: стики, крутилки \sw{S1}/\sw{S2},
        переключатели, триммеры, каналы, логические переключатели, телеметрия, таймеры,
        глобальные переменные, константа \sw{MAX}, гироскоп (если есть) и т.\,д.
  \item[Входы (Inputs)] «пилотская» часть. Здесь стик превращается в управляющую команду:
        с нужным расходом (dual rate), экспонентой, триммом. Входы не знают ничего о модели.
  \item[Миксеры (Mixes)] «конструкторская» часть. Определяют, какие входы попадают в какие
        каналы и как смешиваются: например, элерон и руль высоты на элевонах летающего крыла.
  \item[Выходы (Outputs)] «механическая» часть. Здесь настраивается направление вращения сервы,
        крайние положения (чтобы не упереться в механику) и точная нейтраль.
  \item[Радиомодуль] получает готовые значения каналов и передаёт приёмнику.
\end{description}

Сверху над конвейером находятся «управляющие» механизмы: полётные режимы, логические
переключатели, специальные функции, кривые и глобальные переменные. Они включают и выключают
строки конвейера или меняют их параметры.

\begin{note}
Значения во всём конвейере измеряются в процентах: от $-100\,\%$ до $+100\,\%$ (внутри ---
от $-1024$ до $+1024$). На выходе пульта $\pm100\,\%$ соответствуют длительности импульса
$1500\pm512$\,мкс, т.\,е. примерно 988--2012\,мкс. Включив \menu{Extended limits}
(расширенные пределы) в настройках модели, можно использовать до $\pm150\,\%$.
\end{note}

\section{Главный экран}

После включения (и прохождения предупреждений, см.\ главу~\ref{ch:radio}) вы видите главный экран:

\begin{center}
\lcd{\textbf{Quad-5in}\hfill 7.9V ▮▮▮▯\\
TX 250mW\hfill 14:22\\[2pt]
\hspace*{8mm}\textbf{\large 02:37}\hfill \inv{FM0}\\[2pt]
━ ━ ━ ━ ━\hfill ━ ━ ━ ━ ━\\
SA↑ SB- SC↓ SD↑\hfill RSSI ▮▮▮▮}
\end{center}

На нём отображаются имя модели, напряжение аккумулятора пульта, таймер, текущий полётный режим,
положения триммеров и переключателей. Колесом (или \btn{PAGE>}) можно перелистывать варианты
главного экрана: с положениями стиков, с каналами, с логическими переключателями.

\section{Кнопки и навигация}

На пультах с монохромным экраном управление одинаковое:

\begin{center}\small
\renewcommand{\arraystretch}{1.2}
\begin{tabularx}{\linewidth}{@{}p{33mm}X@{}}
\toprule
Кнопка & Действие \\
\midrule
\btn{SYS} & системное меню (настройки пульта, общие для всех моделей) \\
\btn{MDL} & меню модели (настройки текущей модели) \\
\btn{TELE} & экраны телеметрии текущей модели \\
\btn{PAGE>} / \btn{PAGE<} & следующая / предыдущая страница меню \\
\btn{RTN} & назад, выход, отмена редактирования \\
колесо & перемещение по пунктам, изменение значения \\
\btn{ENTER} (нажатие колеса) & войти в пункт, начать/закончить редактирование \\
долгое \btn{ENTER} & контекстное меню строки (вставить, копировать, переместить, удалить …) \\
долгое \btn{RTN} & быстро вернуться на главный экран (зависит от версии) \\
\bottomrule
\end{tabularx}
\end{center}

\subsection{Редактирование значений}

\begin{enumerate}
  \item Колесом подведите курсор к полю --- оно подсветится инверсией.
  \item Нажмите \btn{ENTER} --- поле начнёт мигать: вы в режиме редактирования.
  \item Вращайте колесо, чтобы изменить значение. Нажмите \btn{ENTER}, чтобы подтвердить,
        или \btn{RTN}, чтобы выйти.
\end{enumerate}

Полезные приёмы:
\begin{itemize}
  \item \textbf{Выбор источника движением.} Когда редактируется поле «источник» или
        «переключатель», просто пошевелите нужным стиком, крутилкой или тумблером --- EdgeTX
        подставит его сам. Так быстрее и без ошибок.
  \item \textbf{Долгое \btn{ENTER} на числовом поле} открывает меню: сбросить к значению
        по умолчанию, инвертировать, задать минимум/максимум или подставить глобальную
        переменную \sw{GV} вместо числа.
  \item \textbf{Долгое \btn{ENTER} на строке списка} (входы, миксеры, логические переключатели)
        --- вставить строку до/после, скопировать, переместить, удалить.
  \item Переключатель в поле \menu{Switch} можно выбрать с «!»: \sw{!SA↑} означает
        «\emph{не} в верхнем положении».
\end{itemize}

\section{Карта меню}

\begin{figure}[H]
\centering
\begin{tikzpicture}[
  every node/.style={font=\sffamily\scriptsize},
  box/.style={draw,rounded corners=2pt,fill=#1,inner sep=4pt,align=left,text width=40mm,anchor=north},
]
\node[box=blkrf,text width=14mm,align=center,anchor=center] (root) at (0,0) {EdgeTX};
\node[box=blkin] (sys) at (-5,-1.1) {\textbf{SYS} --- пульт\\[2pt]
Tools --- Lua и утилиты\\
SD card --- файлы\\
Radio Setup --- общие\\
Global Functions\\
Trainer\\
Hardware --- калибровка\\
About --- версия};
\node[box=blkmix] (mdl) at (0,-1.1) {\textbf{MDL} --- модель\\[2pt]
Model Select --- список\\
Setup --- имя, таймеры, RF\\
Heli --- автомат перекоса\\
Flight Modes\\
Inputs\\
Mixes\\
Outputs\\
Curves\\
Global Variables\\
Logical Switches\\
Special Functions\\
Custom Scripts\\
Telemetry};
\node[box=blkout] (tele) at (5,-1.1) {\textbf{TELE} --- телеметрия\\[2pt]
экраны 1--4:\\
числа, шкалы,\\
Lua-экраны};
\draw[arr] (root.south) -- (mdl.north);
\draw[arr] (root.west) -| (sys.north);
\draw[arr] (root.east) -| (tele.north);
\end{tikzpicture}
\caption{Основные разделы меню EdgeTX на пультах с экраном 128×64}
\end{figure}

\begin{note}
Страница \menu{Heli} появляется, только если прошивка собрана с поддержкой вертолётов
(вариант по умолчанию), а страница \menu{Custom Scripts} --- только в сборках с Lua.
Некоторые страницы можно скрыть в \mpath{SYS\sep Radio Setup} или в настройках модели.
\end{note}

\section{Где хранятся настройки}

\begin{itemize}
  \item \textbf{Настройки пульта} (язык, режим стиков, калибровка, звук) --- файл
        \code{RADIO/radio.yml} на SD-карте.
  \item \textbf{Модели} --- каждая в своём файле \code{MODELS/modelNN.yml}.
  \item Звуки, скрипты, картинки, логи --- в соответствующих папках SD-карты (приложение~\ref{app:sd}).
\end{itemize}

Файлы в формате YAML --- текстовые, их можно открыть на компьютере, сравнить, сохранить
в систему контроля версий или отправить товарищу. Резервная копия SD-карты = резервная копия
всего пульта.

\begin{summary}
\begin{itemize}
  \item EdgeTX --- открытая прошивка, наследница OpenTX.
  \item Сигнал проходит путь: источник → вход → миксер → выход → радиомодуль.
  \item \btn{SYS} --- настройки пульта, \btn{MDL} --- настройки модели, \btn{TELE} --- телеметрия.
  \item Все настройки лежат на SD-карте в YAML-файлах.
\end{itemize}
\end{summary}
